\chapter{Pascal's Triangle}
\label{ch:1}
\section{Pascal's triangle and binomial identities}
\label{sec:1-1}
% TODO: Develop this section; add verified references, examples, and exercises.

\section{Krawtchouk polynomials}
\label{sec:1-2}
For $0\leq n,j\leq N$, define the symmetric Krawtchouk polynomial values by
\begin{equation}\label{eq:gen}
(1+v)^{N-j}(1-v)^j=\sum_{n=0}^N K_n(j;N)v^n.
\end{equation}
Hence
\[
K_n(j;N)=\sum_r(-1)^r\binom jr\binom{N-j}{n-r}.
\]
Binomial coefficients outside their usual range are understood to vanish.
In centered coordinates $x=N-2j$, write
\[
(1+v)^{(N+x)/2}(1-v)^{(N-x)/2}
=\sum_{n=0}^N\frac{v^n}{n!}k_n(x,N).
\]
Thus $k_n(N-2j,N)=n!K_n(j;N)$.

\section{Generating functions}
\label{sec:1-3}
% TODO: Develop this section; add verified references, examples, and exercises.

\section{Krawtchouk matrices}
\label{sec:1-4}
Set $\Phi_N=(K_n(j;N))_{n,j=0}^N$ and
$B_N=\operatorname{diag}(\binom N0,\ldots,\binom NN)$. Binomial duality is
\[
\binom Nj K_n(j;N)=\binom Nn K_j(n;N).
\]
Consequently $S_N=\Phi_NB_N$ is symmetric. The symmetric transform satisfies
\[
\Phi_N^2=2^NI.
\]
For orientation,
\[
\Phi_1=\begin{pmatrix}1&1\\1&-1\end{pmatrix},\qquad
\Phi_2=\begin{pmatrix}1&1&1\\2&0&-2\\1&-1&1\end{pmatrix}.
\]

\section{First view and cross properties}
\label{sec:1-5}
% TODO: Develop this section; add verified references, examples, and exercises.

